\documentclass[11pt,a4paper]{article}

\usepackage[T1]{fontenc}
\usepackage[utf8]{inputenc}
\usepackage{mathpazo}
\usepackage[margin=1in]{geometry}
\usepackage{microtype}
\usepackage{xcolor}
\usepackage{listings}
\usepackage{booktabs}
\usepackage{enumitem}
\usepackage[hidelinks,colorlinks=true,linkcolor=linkblue,urlcolor=linkblue]{hyperref}

\definecolor{linkblue}{HTML}{1F4E79}
\definecolor{codebg}{HTML}{F6F8FA}
\definecolor{codeframe}{HTML}{D6DBE1}
\definecolor{kw}{HTML}{A3254A}
\definecolor{cmt}{HTML}{5A7D5A}
\definecolor{accent}{HTML}{1F4E79}

\lstdefinestyle{sqlstyle}{
  language=SQL,
  morekeywords={ASOF,LATERAL,OVER,RETURNING,EXPLAIN,ANALYZE},
  basicstyle=\ttfamily\footnotesize,
  keywordstyle=\color{kw}\bfseries,
  commentstyle=\color{cmt}\itshape,
  backgroundcolor=\color{codebg},
  frame=single, rulecolor=\color{codeframe},
  framesep=4pt, xleftmargin=5pt,
  breaklines=true, showstringspaces=false,
  columns=fullflexible, keepspaces=true,
}
\lstdefinestyle{cstyle}{
  language=C,
  morekeywords={bool,List,Node,JoinExpr,SelectStmt,RangeSubselect,ColumnRef,A_Expr,SortBy,
                A_Const,ParseState,ListCell},
  basicstyle=\ttfamily\footnotesize,
  keywordstyle=\color{kw}\bfseries,
  commentstyle=\color{cmt}\itshape,
  backgroundcolor=\color{codebg},
  frame=single, rulecolor=\color{codeframe},
  framesep=4pt, xleftmargin=5pt,
  breaklines=true, showstringspaces=false,
  columns=fullflexible, keepspaces=true,
}
\lstset{style=sqlstyle}
\emergencystretch=3em
\hbadness=4000

\newcommand{\code}[1]{\texttt{\small #1}}
\newcommand{\file}[1]{\texttt{\footnotesize #1}}
\newcommand{\asof}{{\normalfont\scshape asof}}

\setlength{\parskip}{0.4em}
\setlength{\parindent}{0pt}

\title{\vspace{-2cm}\textbf{How ASOF JOIN Is Implemented}\\[0.2cm]
       \large Phase 1: rewriting to a LATERAL join}
\author{CS631 — ASOF JOIN project}
\date{\today}

\begin{document}
\maketitle
\vspace{-0.8cm}

\section{What the rewrite does}

\asof{} join is implemented without touching the planner or the executor. The parser
recognises the syntax, and the analyser rewrites the query into an equivalent
\code{LATERAL} join that PostgreSQL already knows how to run. Everything downstream is
untouched code.

\begin{lstlisting}
-- written by the user
SELECT t.sym, t.ts, q.bid
FROM trades t ASOF LEFT JOIN quotes q
  ON t.sym = q.sym AND t.ts >= q.ts;
\end{lstlisting}

\begin{lstlisting}
-- what the database actually analyses
SELECT t.sym, t.ts, q.bid
FROM trades t LEFT JOIN LATERAL (
    SELECT * FROM quotes
    WHERE quotes.sym = t.sym AND t.ts >= quotes.ts
    ORDER BY quotes.ts DESC
    LIMIT 1
) q ON true;
\end{lstlisting}

\code{LATERAL} is what makes the subquery legal: it allows the subquery to reference
\code{t}, the item to its left in the \code{FROM} clause. Semantically the subquery is
evaluated once per row of \code{trades}, with that row's values substituted in, and
\code{ORDER BY ... LIMIT 1} picks the single nearest match. \code{LEFT JOIN ... ON true}
supplies the NULL-extension when no match exists; a plain \code{JOIN} drops the row
instead, which is how \code{ASOF JOIN} and \code{ASOF LEFT JOIN} differ.

\section{Where it happens, and why that choice keeps it small}

PostgreSQL processes a query in stages: \textbf{parse} (text to a raw tree, no catalog
lookups), \textbf{analyse} (resolve names and types), \textbf{rewrite} (views and rules),
\textbf{plan}, \textbf{execute}.

The rewrite is performed at the \emph{start} of analysis, in
\file{src/backend/parser/parse\_clause.c}. \code{transformFromClauseItem()} walks the
\code{FROM} clause; when it meets a \code{JoinExpr} with \code{isAsof} set, it calls
\code{transformAsofToLateralJoin()} and continues with the returned tree:

\begin{lstlisting}[style=cstyle]
/* ASOF joins are executed as LATERAL joins; rewrite before analysis */
if (j->isAsof)
    j = transformAsofToLateralJoin(pstate, j);
\end{lstlisting}

The timing matters. Because this runs \emph{before} names are resolved, the function
manipulates the raw parse tree — \code{ColumnRef} nodes holding strings, not \code{Var}
nodes with resolved table and column numbers. It produces a tree identical to what the
parser would have built had the user typed the \code{LATERAL} form by hand. Name
resolution, type resolution, \code{LATERAL} validity checking, lateral-reference marking
and planning then all happen afterwards in existing code that has no idea \asof{} exists.

That is why Phase 1 is roughly 250 lines rather than several thousand: it is a
source-to-source translation, the technique a compiler calls \emph{desugaring}, where a new
construct is defined entirely by what it translates into.

\section{Building the replacement tree}

\subsection*{Finding the inequality}

The grammar hands over the whole \code{ON} clause as one expression. The function flattens
its top-level \code{AND} into a list of conjuncts and looks for exactly one that is an
inequality relating the two sides:

\begin{itemize}[nosep,leftmargin=*]
\item the operator is one of \code{<}, \code{<=}, \code{>}, \code{>=};
\item one operand is a plain column of the right-hand table;
\item the other references columns, but none belonging to the right-hand table — so it can
      only be evaluated against the left.
\end{itemize}

Finding two such conjuncts, or none, is an error. The remaining conjuncts — the equality
keys and any other filters — are left alone.

\subsection*{Choosing the sort direction}

The operator decides which way "nearest" runs, and therefore the \code{ORDER BY}:

\begin{center}
\begin{tabular}{@{}llll@{}}
\toprule
\textbf{Condition} & \textbf{Meaning} & \textbf{Sort} \\
\midrule
\code{t.ts >= q.ts} & latest at or before   & \code{q.ts DESC} \\
\code{t.ts > q.ts}  & latest strictly before & \code{q.ts DESC} \\
\code{t.ts <= q.ts} & earliest at or after  & \code{q.ts ASC} \\
\code{t.ts < q.ts}  & earliest strictly after & \code{q.ts ASC} \\
\bottomrule
\end{tabular}
\end{center}

Writing the same condition with the operands swapped (\code{q.ts <= t.ts}) means the same
thing, so the code checks both arrangements and flips the direction accordingly.

\subsection*{Assembling it}

The replacement is built from ordinary raw-parse-tree nodes:

\begin{itemize}[nosep,leftmargin=*]
\item a \code{SelectStmt} whose target list is \code{*}, whose \code{FROM} is the original
      right-hand side, whose \code{WHERE} is the \emph{entire original \code{ON} clause},
      with the \code{ORDER BY} above and \code{LIMIT 1};
\item a \code{RangeSubselect} wrapping it with \code{lateral = true}, aliased to the name
      the right table was already visible under — so references like \code{q.bid} elsewhere
      in the query keep resolving;
\item a new \code{JoinExpr} with the original join type, the original left side, that
      subquery as its right side, and \code{ON true}.
\end{itemize}

The input tree is not modified; a new \code{JoinExpr} is returned, because callers may reuse
raw parse trees.

\section{The one subtlety: where the equality keys go}

The equality keys go \emph{inside} the subquery's \code{WHERE}, together with the
inequality. Putting them in the outer \code{ON} instead looks equivalent and is not:

\begin{lstlisting}
-- WRONG
trades t LEFT JOIN LATERAL (
    SELECT * FROM quotes WHERE quotes.ts <= t.ts
    ORDER BY quotes.ts DESC LIMIT 1
) q ON t.sym = q.sym
\end{lstlisting}

That subquery finds the latest quote \emph{of any symbol} before the trade, and the outer
condition then rejects it if the symbol happens to differ — returning NULL even though a
perfectly good quote for the right symbol exists. Filtering and "take the latest" do not
commute: you must restrict the candidates \emph{before} choosing the best one, never after.

This is also why PostgreSQL's optimiser could never perform this rewrite for us. Pushing a
predicate below a \code{LIMIT} changes the result, so the planner correctly refuses to do
it; the predicate has to be placed correctly at construction time.

\section{What it costs, and what it does not do}

The resulting plan is a nested loop with the subquery on the inner side:

\begin{lstlisting}[style=sqlstyle]
Nested Loop Left Join
  ->  Seq Scan on trades t
  ->  Limit
        ->  Index Scan Backward using quotes_sym_ts_idx on quotes q
              Index Cond: ((sym = t.sym) AND (ts <= t.ts))
\end{lstlisting}

With a suitable index this is one index descent per left row: $O(n \log m)$, with random
access per probe. Without one it degrades to a sort or scan per left row, which is
$O(n \cdot m)$. Phase 1 therefore makes \asof{} \emph{expressible}, not fast — it performs
exactly as well as the hand-written \code{LATERAL} query, because it is that query.

Making it fast is Phase 2: sort both inputs by \code{(key, time)} once and walk them
together, keeping a single remembered row, which is $O(n + m)$ — and free of sorting
entirely when an index already supplies the order. That plan cannot be expressed as a
\code{LATERAL} join at any price, which is the reason to build a real operator.

\subsection*{Known limitations of the current implementation}

\begin{itemize}[nosep,leftmargin=*]
\item Views store the rewritten \code{LATERAL} form, so \code{pg\_get\_viewdef} does not
      show \asof{} syntax.
\item Operands are classified by matching column-reference names textually, since the
      rewrite runs before name resolution; an unqualified reference is never attributed to
      the right-hand table.
\item The inequality operator is not checked for a usable sort order — which Phase 2 needs.
\item The right-hand side must be a table or an aliased subquery.
\item The right operand of the inequality must be a bare column, so
      \code{ON t.ts >= q.ts + interval '1 min'} is rejected.
\item Ties on the match column resolve arbitrarily, as \code{LIMIT 1} does.
\end{itemize}

\end{document}
